\documentclass[11pt,a4paper]{article}

\usepackage[T1]{fontenc}
\usepackage[utf8]{inputenc}
\usepackage[french]{babel}
\usepackage{lmodern}
\usepackage{geometry}
\usepackage{amsmath,amssymb}
\usepackage{booktabs}
\usepackage{graphicx}
\usepackage{float}
\usepackage{hyperref}
\usepackage{xcolor}
\geometry{margin=2.5cm}

\title{Quantum classification on a few qubit}
\author{Nom Prenom\\2A -- Progra Quantique}
\date{\today}

\begin{document}
\maketitle
\tableofcontents
\newpage

\section{Introduction}
La classification supervisee est un probleme central en apprentissage automatique. Dans ce projet, on etudie une question simple mais interessante: \emph{jusqu'ou peut-on aller avec un seul qubit pour classer des donnees a plusieurs features}?

L'idee est de construire un classifieur quantique leger, puis de compenser ses limites de capacite avec un schema d'ensemble. Le fil directeur est:
\begin{itemize}
  \item un classifieur QNN a 1 qubit;
  \item une validation sur le probleme du cercle (frontiere non lineaire);
  \item une amelioration des performances via AdaBoost.
\end{itemize}

\section{Idee de base: un qubit pour classer des donnees}
\subsection{Representation des donnees}
On considere des entrees $x \in \mathbb{R}^n$. Dans l'experience principale, le cas est bidimensionnel ($n=2$), mais l'idee generale consiste a reutiliser un petit espace quantique (ici un seul qubit) avec une profondeur de circuit suffisante et des parametres entrainables.

Dans le code, chaque couche applique:
\begin{enumerate}
  \item une porte de melange (Hadamard);
  \item une transformation parametrique $U(\phi_0,\phi_1,\phi_2)$, ou les angles dependent des features et de parametres libres.
\end{enumerate}

Une forme type est:
\[
\phi_k = \theta_k + \omega_k \cdot x_k,
\]
avec adaptation pratique pour le cas 2D (ajout d'une troisieme composante nulle dans l'implementation).

\subsection{Decision de classification}
Le circuit retourne des probabilites $(p_0, p_1)$ par mesure dans la base computationnelle. La regle de decision binaire est:
\[
\hat y =
\begin{cases}
0 & \text{si } p_0 \ge 0.5,\\
1 & \text{sinon.}
\end{cases}
\]

\section{Cas d'etude: le probleme du cercle}
Le dataset est genere dans $[-1,1]^2$. L'etiquette est definie par:
\[
y =
\begin{cases}
0 & \text{si } x^2 + y^2 < \frac{2}{\pi},\\
1 & \text{sinon.}
\end{cases}
\]

Ce probleme est pertinent car la frontiere de decision est non lineaire (cercle), ce qui teste la capacite du modele a representer des formes non affines.

\section{Entrainement du modele QNN simple}
\subsection{Fonction de cout et optimisation}
Le modele simple est optimise avec COBYLA, sans gradient analytique, via une fonction de cout quadratique entre probabilites predites et cibles one-hot.

Forme generique:
\[
\mathcal{L}(\Theta) = \frac{1}{N}\sum_{i=1}^N \left[(p_0(x_i)-e_{0,i})^2 + (p_1(x_i)-e_{1,i})^2\right],
\]
ou $(e_{0,i}, e_{1,i})$ est la cible one-hot.

\subsection{Simulation avec et sans bruit}
Deux regimes sont utilises:
\begin{itemize}
  \item \textbf{Noiseless}: simulation exacte (statevector);
  \item \textbf{Noisy}: simulation a shots avec erreur depolarisante.
\end{itemize}

\section{Extension par AdaBoost}
\subsection{Pourquoi AdaBoost ici}
Un classifieur a 1 qubit peut rester limite sur certaines zones de l'espace. AdaBoost permet de combiner plusieurs \emph{weak learners} pour former un classifieur global plus robuste.

\subsection{Algorithme implemente}
On initialise des poids uniformes sur les echantillons:
\[
w_i^{(1)} = \frac{1}{N}.
\]

Pour chaque round $t=1,\dots,T$:
\begin{enumerate}
  \item entrainer un weak learner QNN avec la loss \emph{ponderee} par $w_i^{(t)}$;
  \item calculer l'erreur ponderee:
  \[
  \epsilon_t = \sum_i w_i^{(t)}\,\mathbf{1}[h_t(x_i)\neq y_i];
  \]
  \item calculer le poids du learner:
  \[
  \alpha_t = \frac{1}{2}\ln\left(\frac{1-\epsilon_t}{\epsilon_t}\right);
  \]
  \item mettre a jour les poids des exemples:
  \[
  w_i^{(t+1)} \leftarrow w_i^{(t)}\exp\left(-\alpha_t y_i h_t(x_i)\right),
  \]
  puis normaliser.
\end{enumerate}

La prediction finale est un vote pondere:
\[
F(x) = \sum_{t=1}^T \alpha_t h_t(x),
\qquad
\hat y = \mathbf{1}[F(x)\ge 0].
\]

\section{Resultats experimentaux}
\subsection{Metriques}
Les metriques utilisees sont:
\begin{itemize}
  \item accuracy,
  \item precision,
  \item recall,
  \item matrice de confusion.
\end{itemize}

\subsection{Tableau de comparaison}
Remplacer les valeurs ci-dessous par les valeurs du notebook.

\begin{table}[H]
\centering
\begin{tabular}{lccc}
\toprule
Modele & Accuracy & Precision & Recall \\
\midrule
Single QNN & \texttt{a completer} & \texttt{a completer} & \texttt{a completer} \\
AdaBoost-QNN & \texttt{a completer} & \texttt{a completer} & \texttt{a completer} \\
\bottomrule
\end{tabular}
\caption{Comparaison des performances sur le jeu de test.}
\end{table}

\subsection{Matrices de confusion}
Tu peux inserer les figures produites dans le notebook.

\begin{figure}[H]
  \centering
  % \includegraphics[width=0.8\textwidth]{RESULTS/confusion_matrices.png}
  \fbox{\parbox[c][5cm][c]{0.85\textwidth}{\centering Figure des matrices de confusion a inserer.}}
  \caption{Matrices de confusion: Single QNN vs AdaBoost-QNN.}
\end{figure}

\section{Discussion}
\begin{itemize}
  \item Le modele 1 qubit montre qu'une architecture tres compacte peut deja capturer une frontiere non lineaire.
  \item La limitation principale est la capacite de representation et la sensibilite a l'optimisation locale.
  \item AdaBoost ameliore en pratique les performances en reorientant l'apprentissage vers les exemples difficiles.
\end{itemize}


\appendix
\section{Hyperparametres (a renseigner)}
\begin{itemize}
  \item Nombre de points: $N=\,$\texttt{...}
  \item Fraction test: \texttt{...}
  \item Profondeur du circuit ($RC$): \texttt{...}
  \item Iterations max COBYLA: \texttt{...}
  \item Nombre de weak learners AdaBoost: \texttt{...}
  \item Regime bruit/noiseless: \texttt{...}
\end{itemize}

\end{document}
